\begin{lemma}
Let \(\eta,\varepsilon_1,\varepsilon_2\) be positive real constants satisfying
\[
0<\varepsilon_2<\varepsilon_1<\eta<1,\qquad \eta<\frac1{16},
\]
\[
3\varepsilon_2\le\frac{\varepsilon_1}{10},
\qquad
8\sqrt{\varepsilon_1}+10\varepsilon_1\le\frac{\eta^3}{2}.
\]
Then there is \(n_0\) such that
\[
\varepsilon_2^{-2}\le n_0
\]
and
\[
\noderef{ParameterHierarchyEventualComparisons}(\eta,\varepsilon_1,
\varepsilon_2,n_0)
\]
holds.
\end{lemma}
\begin{proof}
Choose \(n_0\) at least \(\varepsilon_2^{-2}\), and then enlarge it finitely
many times.  For \(n>n_0\), write
\[
L=\log n,\quad
k_R=(1+\eta)\sqrt{nL},\quad
K=\noderef{TwoBiteNaturalIndependenceScale}(1+\eta,n),\quad k=(K:\mathbb R),
\]
\[
m_R=\frac{n}{L^2},\quad
M=\noderef{TwoBiteNaturalReducedVertexCount}(n),\quad m=(M:\mathbb R),
\]
\[
p=\frac12\sqrt{\frac Ln},\quad
t_1=\frac{\sqrt{nL}}{\log L},\quad
t_2=n^{1/4+\eta},\quad t_3=n^{2\eta}.
\]
First enlarge \(n_0\) so that, for every \(n>n_0\),
\[
n>1,\quad L>1,\quad \log L>0,\quad k_R>2,\quad m_R>2,
\quad \sqrt{L/n}\le1,
\]
\[
\frac{k_R+1}{n}\le1,\qquad
\frac{\sqrt n}{2L^{3/2}}\ge1.
\]
These are eventual conditions because \(L\to\infty\), \(\log L\to\infty\),
\(\sqrt{nL}\to\infty\), \(n/L^2\to\infty\), \(\sqrt{L/n}\to0\),
\((k_R+1)/n\to0\), and \(\sqrt n/(2L^{3/2})\to\infty\).
The node \(\noderef{ParameterHierarchyBaseThreshold}\) packages this
base-threshold choice and the initial estimates derived from it.

Since \(L>1\), both \(k_R\) and \(m_R\) are nonnegative.  Applying
\(\noderef{NaturalParameterApproximation}\) with \(\kappa=1+\eta\) gives
\[
k_R\le k<k_R+1,\qquad m\le m_R<m+1.
\]
The bounds \(k_R>2\) and \((k_R+1)/n\le1\) imply \(k<2k_R\) and \(K\le n\).
The bound \(m_R>2\) implies
\[
\frac12m_R\le m\le m_R,\qquad 0<m.
\]
Moreover \(p\ge0\), \(p\le1/2\), and
\[
2pm\ge p m_R
=\frac12\sqrt{\frac Ln}\frac n{L^2}
=\frac{\sqrt n}{2L^{3/2}}\ge1.
\]
This proves the first eight positivity and rounding conjuncts in
\(\noderef{ParameterHierarchyEventualComparisons}\).

For \(k>2\), applying \(\noderef{RealChooseTwoQuadraticBounds}\) to \(k\)
gives
\[
\frac{k^2}{4}\le \noderef{RealChooseTwo}(k)\le\frac{k^2}{2}. \tag{C}
\]
For every \(x\ge0\), the same helper also gives
\[
\noderef{RealChooseTwo}(x)\le x^2. \tag{U}
\]
After enlarging the threshold so that \(\log L\ge1\), the estimate
\[
\frac{2k}{t_1}+1<4(1+\eta)\log L+1\le5(1+\eta)\log L \tag{D}
\]
follows from \(k<2k_R\) and \(t_1=\sqrt{nL}/\log L\).

We now enlarge \(n_0\) to make each of the following displayed threshold
conditions hold for every \(n>n_0\).  Each left side tends to \(0\), because
it is a fixed constant times a product of powers of \(L=\log n\) and
\(\log L=\log\log n\) divided by a positive power of \(n\); each displayed
lower-growth condition tends to infinity.  The exponents of \(n\) are positive
because \(0<\eta<1/16\).

For
\[
n^{4\eta}k\le(\varepsilon_1/2)k^2
\]
it is enough, using \(k\ge k_R\), that
\[
\frac{n^{4\eta}}{k_R}
=\frac{1}{1+\eta}\frac{n^{-1/2+4\eta}}{\sqrt L}
\le\frac{\varepsilon_1}{2}. \tag{T1}
\]
For
\[
2kL^2\le(\varepsilon_1/8)k^2
\]
it is enough that
\[
\frac{2L^2}{k_R}
=\frac{2}{1+\eta}\frac{L^{3/2}}{\sqrt n}
\le\frac{\varepsilon_1}{8}. \tag{T2}
\]
The three-term comparison
\[
\frac1{2L}\noderef{RealChooseTwo}(k)+2kL^2+
\frac1{\sqrt L}\noderef{RealChooseTwo}(k)
\le\frac2{\sqrt L}\noderef{RealChooseTwo}(k)
\]
follows from (C), \(L>1\), and
\[
\frac{16L^{5/2}}{k_R}\le1. \tag{T3}
\]
For the exponential binomial inequality,
\(\noderef{ParameterHierarchyT4ChooseExpBound}\) gives
\(\binom nK\le\exp(kL)\).  The threshold
\(\noderef{ParameterHierarchyT4ExponentThreshold}\), together with
\(\noderef{ParameterHierarchyT4ExponentAlgebra}\), gives
\[
2kL+L\le
\frac{\noderef{RealChooseTwo}(k)^2}{Lm n^{8\eta}}.
\]
Applying \(\noderef{ParameterHierarchyT4ExponentialConclusion}\) to these
two estimates yields
\[
\exp\!\left(
-\frac{\noderef{RealChooseTwo}(k)^2}{Lm n^{8\eta}}
\right)\binom nK\le n^{-1}. \tag{T4}
\]

For
\[
\frac2{\sqrt L}\noderef{RealChooseTwo}(k)
\le\frac{\varepsilon_1}{2}k^2,
\]
the threshold \(\noderef{ParameterHierarchyT5SqrtLogThreshold}\) gives,
after enlarging \(n_0\),
\[
\frac1{\sqrt L}\le\frac{\varepsilon_1}{2}. \tag{T5}
\]
Then \(\noderef{ParameterHierarchyT5ChooseBound}\), applied with (C) and
the positivity facts \(0\le k\), \(2\le k\), and \(\sqrt L>0\), gives the
displayed inequality.
For
\[
4\log k\le\frac{\eta^3}{2}pk^2,
\]
the threshold \(\noderef{ParameterHierarchyT6LogThreshold}\) gives
\[
\frac{8\log(2(1+\eta)\sqrt{nL})}
{(1+\eta)^2\sqrt n\,L^{3/2}}
\le\frac{\eta^3}{2}. \tag{T6}
\]
Together with \(k<2k_R\), \(k_R\le k\), and
\[
pk_R^2=\frac{(1+\eta)^2}{2}\sqrt n\,L^{3/2}.
\]
the algebra helper \(\noderef{ParameterHierarchyT6LogAlgebra}\) gives the
displayed inequality.
For
\[
t_1k\le\varepsilon_1 k^2
\]
the threshold \(\noderef{ParameterHierarchyT7LogLogThreshold}\) gives,
after enlarging \(n_0\),
\[
\frac{t_1}{k_R}=\frac1{(1+\eta)\log L}\le\varepsilon_1. \tag{T7}
\]
Together with \(k_R\le k\), the algebra helper
\(\noderef{ParameterHierarchyT7Algebra}\) gives the displayed inequality.
Using (D) and (U), the inequality
\[
\noderef{RealChooseTwo}\left(\frac{2k}{t_1}+1\right)100L^3
\le\frac12k
\]
follows from the threshold \(\noderef{ParameterHierarchyT8Threshold}\),
which after enlarging \(n_0\) gives
\[
\frac{200(5(1+\eta))^2(\log L)^2L^3}
{(1+\eta)\sqrt{nL}}\le1. \tag{T8}
\]
Together with \(k_R\le k\), \(0\le k\), \(0<t_1\), and (D), the algebra
helper \(\noderef{ParameterHierarchyT8Algebra}\) gives the displayed
inequality.
For
\[
\frac{2k}{t_1}(2pm)\le\frac{\varepsilon_1}{10}k,
\]
the estimate \(\noderef{ParameterHierarchyT9MassBound}\) gives
\[
2pm\le\frac{\sqrt n}{L^{3/2}}.
\]
The threshold \(\noderef{ParameterHierarchyT9LogThreshold}\), after
enlarging \(n_0\), gives
\[
\frac{2\log L}{L^2}\le\frac{\varepsilon_1}{10}. \tag{T9}
\]
Together with \(t_1=\sqrt{nL}/\log L\) and \(0\le k\), the algebra helper
\(\noderef{ParameterHierarchyT9Algebra}\) gives the displayed inequality.
For
\[
\frac{2k}{t_1}\noderef{RealChooseTwo}(2pm)\le\varepsilon_1 k^2,
\]
use (U) and the mass estimate
\(\noderef{ParameterHierarchyT9MassBound}\),
\[
2pm\le\frac{\sqrt n}{L^{3/2}}.
\]
The threshold \(\noderef{ParameterHierarchyT10LogThreshold}\), after
enlarging \(n_0\), gives
\[
\frac{2\log L}{(1+\eta)L^4}\le\varepsilon_1. \tag{T10}
\]
Together with \(t_1=\sqrt{nL}/\log L\), \(k_R\le k\), \(0\le k\), and
\(0\le2pm\), the algebra helper
\(\noderef{ParameterHierarchyT10Algebra}\) gives the displayed inequality.
For
\[
\noderef{RealChooseTwo}\left(\frac{2k}{t_1}+1\right)200L^3
\le\frac{\varepsilon_1}{10}k
\]
the threshold \(\noderef{ParameterHierarchyT11Threshold}\), after enlarging
\(n_0\), gives
\[
\frac{2000(5(1+\eta))^2(\log L)^2L^3}
{\varepsilon_1(1+\eta)\sqrt{nL}}\le1. \tag{T11}
\]
Together with (D), \(0\le k\), \(0<t_1\), and \(k_R\le k\), the algebra
helper \(\noderef{ParameterHierarchyT11Algebra}\) gives the displayed
inequality.
For
\[
\noderef{RealChooseTwo}\left(\frac{2k}{t_1}+1\right)
\noderef{RealChooseTwo}(200L^3)\le\frac{\varepsilon_1}{10}k^2
\]
the threshold \(\noderef{ParameterHierarchyT12Threshold}\), after
enlarging \(n_0\), gives
\[
\frac{400000(5(1+\eta))^2(\log L)^2L^5}
{\varepsilon_1(1+\eta)^2n}\le1. \tag{T12}
\]
Together with (D), \(1\le L\), \(0\le k\), \(0<t_1\), and \(k_R\le k\),
the algebra helper \(\noderef{ParameterHierarchyT12Algebra}\) gives the
displayed inequality.
The three constant conjuncts
\[
3\varepsilon_2\le\varepsilon_1/10,\qquad
0\le\varepsilon_2,\qquad
\varepsilon_2\le1
\]
come from the assumptions \(3\varepsilon_2\le\varepsilon_1/10\),
\(0<\varepsilon_2\), and \(\varepsilon_2<\varepsilon_1<\eta<1\).

For
\[
2pm\le\frac{\varepsilon_2}{100}t_1
\]
the helper \(\noderef{ParameterHierarchyT16LogThreshold}\), after enlarging
\(n_0\), gives the displayed inequality from the mass estimate
\[
2pm\le\frac{\sqrt n}{L^{3/2}},
\]
the positivity of \(n\), \(L\), and \(\log L\), and the identity
\(t_1=\sqrt{nL}/\log L\).  This helper obtains the needed threshold by
instantiating \(\noderef{ParameterHierarchyT9LogThreshold}\) with
\(\varepsilon_2/5\), so that
\[
\frac{2\log L}{L^2}\le\frac{\varepsilon_2}{50}
\quad\text{and hence}\quad
\frac{\log L}{L^2}\le\frac{\varepsilon_2}{100}. \tag{T16}
\]
We next prove the first \(t_1\)-relative error conjunct,
\[
\left(\frac{2k}{t_1}+1\right)100L^3
\le\frac{\varepsilon_2}{100}t_1.
\]
Put
\[
A=\frac{2k}{t_1}+1,\qquad S=\sqrt{nL},\qquad q=\log L.
\]
The positivity facts above give \(S>0\), \(q>0\), and \(t_1=S/q\).  From
(D) we have
\[
A\le 5(1+\eta)q.
\]
Since \(L>0\), multiplication by \(100L^3\ge0\) gives
\[
A\,100L^3\le 100\cdot5(1+\eta)qL^3. \tag{E17}
\]
After enlarging \(n_0\), the helper
\(\noderef{ParameterHierarchyT17T19Threshold}\) supplies the T17, T18, and
T19 thresholds.
The displayed T17 threshold is
\[
\frac{100\cdot5(1+\eta)(\log L)^2L^3}{\sqrt{nL}}
\le\frac{\varepsilon_2}{100}, \tag{T17}
\]
that is,
\[
\frac{100\cdot5(1+\eta)q^2L^3}{S}
\le\frac{\varepsilon_2}{100}.
\]
Multiplying this inequality by the positive factor \(S/q=t_1\) gives
\[
100\cdot5(1+\eta)qL^3
\le\frac{\varepsilon_2}{100}t_1. \tag{F17}
\]
Combining (E17) and (F17) proves the T17 conjunct.

The same calculation gives the two immediately following \(t_1\)-relative
error conjuncts.  From (D) and \(L>0\),
\[
A\,200L^3\le 200\cdot5(1+\eta)qL^3.
\]
The T18 threshold
\[
\frac{200\cdot5(1+\eta)(\log L)^2L^3}{\sqrt{nL}}
\le\frac{\varepsilon_2}{100}, \tag{T18}
\]
after multiplication by \(S/q=t_1\), gives
\[
200\cdot5(1+\eta)qL^3
\le\frac{\varepsilon_2}{100}t_1.
\]
Thus
\[
\left(\frac{2k}{t_1}+1\right)200L^3
\le\frac{\varepsilon_2}{100}t_1.
\]
Likewise, from (D) and \(L>0\),
\[
A\,400L^5\le 400\cdot5(1+\eta)qL^5.
\]
The T19 threshold
\[
\frac{400\cdot5(1+\eta)(\log L)^2L^5}{\sqrt{nL}}
\le\frac{\varepsilon_2}{100}, \tag{T19}
\]
after multiplication by \(S/q=t_1\), gives
\[
400\cdot5(1+\eta)qL^5
\le\frac{\varepsilon_2}{100}t_1,
\]
and hence
\[
\left(\frac{2k}{t_1}+1\right)400L^5
\le\frac{\varepsilon_2}{100}t_1.
\]
For
\[
t_2n^{1/4-\eta}\le\varepsilon_1 k
\]
unfold \(t_2=n^{1/4+\eta}\), so the left side is \(n^{1/2}\).  Since
\(k\ge(1+\eta)\sqrt n\sqrt L\), it is enough that
\[
\frac1{\varepsilon_1(1+\eta)\sqrt L}\le1. \tag{T20}
\]
For the cutoff comparison
\[
\noderef{TwoBiteLargeCutoff}(\eta,n)<\noderef{TwoBiteHugeCutoff}(n),
\]
the definitions reduce it to
\[
\frac{\log L}{n^{1/4-\eta}\sqrt L}<1. \tag{T21}
\]
For the huge-cutoff upper bound, use \(\varepsilon_2\le1\), \(m\le m_R\),
and \(p=\frac12\sqrt{L/n}\).  The two summands are at most
\(2Lt_2\) and \(\sqrt n/\sqrt L\), so it is enough that
\[
\frac{2(\log L)L^{1/2}}{n^{1/4-\eta}}<\frac12,\qquad
\frac{\log L}{L}<\frac12. \tag{T22}
\]
For
\[
100(n^{1/4}+1/2)L^3\le\noderef{TwoBiteLargeCutoff}(\eta,n),
\]
first require \(n^{1/4}\ge1\); then it is enough that
\[
\frac{200L^3}{n^\eta}\le1. \tag{T23}
\]
For
\[
k<n^{1/4}t_2-\noderef{RealChooseTwo}(n^{1/4})100L^3,
\]
use (U) to bound the subtracted term by \(100n^{1/2}L^3\).  The assertion
follows from
\[
\frac{100L^3}{n^\eta}\le\frac12,\qquad
\frac{2(1+\eta)\sqrt L}{n^\eta}<\frac12. \tag{T24}
\]
For
\[
t_3^2k\le\frac{\varepsilon_1}{2}k^2
\]
unfold \(t_3=n^{2\eta}\); this is again (T1). \tag{T25}
For
\[
k\left(\frac{2k}{L}+n^{1/4}\right)\le\varepsilon_1 k^2,
\]
divide by \(k>0\) and require
\[
\frac2L\le\frac{\varepsilon_1}{2},\qquad
\frac1{(1+\eta)n^{1/4}\sqrt L}\le\frac{\varepsilon_1}{2}. \tag{T26}
\]

For the combined coefficient inequality, the assumed bound
\[
8\sqrt{\varepsilon_1}+10\varepsilon_1\le\eta^3/2
\]
multiplied by \(pk^2\ge0\) controls the first two terms by
\((\eta^3/2)pk^2\), and (T6) controls \(4\log k\) by the same quantity.
Adding gives
\[
8\sqrt{\varepsilon_1}pk^2+10\varepsilon_1pk^2+4\log k
\le\eta^3pk^2. \tag{T27}
\]
For
\[
\exp(kL+2kn^{1/4-3\eta}L-n^{3/4-2\eta})\le\varepsilon_1,
\]
use \(k<2(1+\eta)\sqrt n\sqrt L\).  It is enough that
\[
2(1+\eta)n^{1/2}L^{3/2}\le\frac14n^{3/4-2\eta},
\qquad
4(1+\eta)n^{3/4-3\eta}L^{3/2}\le\frac14n^{3/4-2\eta},
\]
and then, after the exponent is at most \(-\frac12n^{3/4-2\eta}\), that
\[
\exp\left(-\frac12n^{3/4-2\eta}\right)\le\varepsilon_1. \tag{T28}
\]
For the last exponential, \(k\ge k_R\) gives
\[
\frac{\eta}{4}kL\ge\frac{\eta(1+\eta)}4\sqrt n\,L^{3/2}.
\]
Thus it is enough that
\[
\exp\left(-\frac{\eta(1+\eta)}4\sqrt n\,L^{3/2}\right)
\le\varepsilon_1. \tag{T29}
\]

The finitely many threshold requirements above can be imposed simultaneously
by replacing \(n_0\) with their maximum.  For every \(n>n_0\), the preceding
paragraphs prove each conjunct in
\(\noderef{ParameterHierarchyEventualComparisons}(\eta,\varepsilon_1,
\varepsilon_2,n_0)\), and the initial choice of \(n_0\) gives
\(\varepsilon_2^{-2}\le n_0\).
\end{proof}
